% paper/section_4_deployment.tex
% REWRITTEN evaluation section for the deterministic-primary architecture.
% Supersedes the cosine-primary/COMPASS/OpenClaw-hookgap evaluation this
% section previously described.

\section{Evaluation}
\label{sec:eval}

We evaluate Kavach on three independent axes: deterministic-primary
detection on a large reference attack corpus (\S\ref{sec:vault-executor-eval}),
CHANNEL's session-level provenance mechanism against real multi-step attack
and benign trajectories (\S\ref{sec:channel-eval}), and a large-scale
real-trajectory validation against AgentDojo's own recorded model
transcripts, not a benchmark replay of minimal ground-truth call sequences
(\S\ref{sec:gpt4o}). Throughout, we report benign false-positive counts
jointly with attack recall, and we distinguish measurements taken via
in-process scoring from measurements taken via the live HTTP dispatch path,
since these two paths previously diverged in ways that mattered
(\S\ref{sec:audit-resolution}).

\subsection{An audit discrepancy, and how it was resolved}
\label{sec:audit-resolution}
Early in this work, three different measurements of VAULT/EXECUTOR detection
rate on the same reference corpus disagreed sharply: 34.2\%/34.5\%
(deterministic layer alone, in-process), 72.6\%/89.9\% (full pipeline,
in-process), and 14.9\%/19.2\% (an earlier live-HTTP measurement). We traced
this to two independent causes rather than accepting any one number at face
value: the live-HTTP figure was corrupted by a duplicate server process
answering a fraction of requests, and, separately, the metric itself
differed --- one measurement counted \textsc{Block}-or-\textsc{Escalate} as
a hit while another counted only the deterministic layer's contribution.
Once both were corrected, the live-HTTP and in-process full-pipeline numbers
agreed. We report this resolution explicitly because a paper that reports
only the final number, without the discrepancy that preceded it, would
understate how easily a benchmark harness's own bugs can masquerade as
detector performance --- a risk equally present in any comparably complex
evaluation pipeline, ours included.

\subsection{VAULT and EXECUTOR: reference-corpus detection}
\label{sec:vault-executor-eval}
On the 519-case AtomicRedTeam/GTFOBins reference corpus (117 VAULT-targeted,
203 EXECUTOR-targeted cases; the corpus's \texttt{target\_patterns} field is
unmodified since its original authoring commit, confirmed via \texttt{git
diff}, so it is independent ground truth rather than tuned after the fact),
the deterministic layer alone recovers VAULT 34.2\% (40/117) and EXECUTOR
34.5\% (70/203); adding the cosine-escalate-only triage layer over the
deterministic residual (\S\ref{sec:triage}) raises full-pipeline detection to
VAULT 74.4\% (87/117) and EXECUTOR 90.1\% (181/201), with zero new hard-block
false positives on either benign population we measure (793 general-purpose
BFCL calls, 55 agent-shaped benign calls). A held-out even/odd split of the
same population confirms the full-pipeline rate is stable rather than an
artifact of the specific partition: VAULT 71.2\% (42/59) on the even half
versus 74.1\% (43/58) on the odd half; EXECUTOR 89.2\% (91/102) versus 91.1\%
(92/101). We report both halves explicitly rather than only their average,
consistent with our standing discipline of never characterizing a split as
``balanced'' or ``generalizing'' from an aggregate number alone.

\subsection{CHANNEL: session-level provenance}
\label{sec:channel-eval}
CHANNEL's provenance mechanism (\S\ref{sec:channel}) is evaluated on two real
populations: 22 real multi-step benign sessions and the AgentDojo attack
cases whose two-call read-then-send shape is structurally applicable to
CHANNEL's mechanism. On the benign population, 5/22 sessions correctly
escalate (0 hard-block false positives); these are exactly the sessions in
which a legitimate destination value (an IBAN, most commonly) was read from
a file or transaction history earlier in the same session rather than typed
by the user, which is precisely the ambiguity that motivated an escalate,
not block, verdict for this mechanism (\S\ref{sec:channel}). On the attack
side, 5/5 structurally applicable AgentDojo cases correctly escalate, and we
verified this twice: once against the benchmark's minimal
\texttt{ground\_truth()} call sequence, and again, independently, against
AgentDojo's own full recorded model trajectories for the same five cases ---
real, longer sessions with additional benign calls interleaved, real tool
outputs, and the actual sequence a deployed model produced rather than the
benchmark's answer key. This second check caught a bug in our own
verification script (it tracked only the final step's verdict, so a benign
call after the real attack step silently overwrote an earlier correct
escalation) before we could report the number as final; the corrected script
confirms genuine 5/5 on the real trajectories.

\subsection{Large-scale real-trajectory validation}
\label{sec:gpt4o}
The measurements above use either a synthetic reference corpus or a small,
hand-verified set of AgentDojo cases. To validate detection at scale on real
model behavior, we replayed all 629 real recorded attack trajectories for
GPT-4o (\texttt{gpt-4o-2024-05-13}) across AgentDojo's four suites'
\texttt{important\_instructions} attack family, sourced directly from
AgentDojo's own published \texttt{runs/} data --- real conversation
transcripts with real tool calls and real tool outputs, not a benchmark
replay of minimal ground-truth sequences.

\paragraph{Headline result.} Kavach fires (\textsc{Block} or
\textsc{Escalate}, any minister, at any point in the session) on 467/629
(74.2\%) of these real trajectories: banking 119/144 (82.6\%), slack 105/105
(100\%), travel 63/140 (45.0\%), workspace 180/240 (75.0\%).

\paragraph{Diagnosing the apparent misses, rather than reporting them flat.}
Of the 162 trajectories on which Kavach did not fire, we traced every case
to its actual tool-call content rather than reporting a bare miss rate.
AgentDojo's own \texttt{security} field, which we cross-reference, marks a
trajectory as a successful attack whenever the environment state changed at
all relative to the pre-trajectory state and does not exactly match the
injected goal --- which conflates ``the injected attack succeeded'' with
``the model made no consequential call related to the injection at all.''
131 of the 162 misses fall in the latter category: the model never invoked
any consequential tool in connection with the injected instruction, so there
was no unauthorized action for any minister to catch. We verified this
directly by extracting every tool call from each miss's recorded
transcript and checking it against a fixed consequential-tool list; a
trajectory whose only calls are read-only lookups (hotel ratings,
restaurant listings, calendar reads) is not a detection gap, regardless of
what AgentDojo's own field reports.

\paragraph{The remaining 19 real misses, categorized by mechanism.} The
other 19 trajectories do contain a genuine unauthorized consequential call
that Kavach did not catch, and we categorize each by whether the underlying
attack has a destination \emph{value} to trace (CHANNEL's domain) or is
purely an unauthorized \emph{action} with no data flow at all (a question
closer to NAVIGATOR's domain, \S\ref{sec:limitations}). Five cases
(\texttt{travel}, all \texttt{reserve\_hotel}) name an attacker-substituted
external party --- a real destination value, of the same shape as
\texttt{send\_money}'s recipient IBAN --- but the tool was never added to
CHANNEL's destination-tool list, a scope gap rather than a mechanism
failure. Fourteen cases (\texttt{banking}: 11 cases,
\texttt{update\_scheduled\_transaction}/\texttt{update\_user\_info};
\texttt{workspace}: 3 cases, \texttt{delete\_file}/\texttt{append\_to\_file})
mutate existing account state or same-account data with no new destination
value introduced at all --- there is no value for any provenance mechanism,
however general, to trace, because the question these attacks pose is
``should this action have happened,'' not ``where did this value come
from.'' We attempted to close the five \texttt{reserve\_hotel} cases by
adding the tool to CHANNEL's destination-tool list, and found this
introduces a new false positive: CHANNEL's destination-value extractor
recognizes only email/IBAN/phone/URL shapes, so a proper-noun business name
such as a hotel is never recognized as ``present in the user's own
instruction'' even when it literally is, causing a legitimate reservation
the user explicitly named to be misclassified \textsc{novel} and incorrectly
escalate. We reverted this change rather than ship a fix with a known new
false-positive class, and report the gap as a documented scope limitation
(\S\ref{sec:limitations}) rather than a closed fix.

\subsection{InjecAgent}
\label{sec:injecagent}
On the full InjecAgent corpus (544 cases), Kavach's deterministic-primary
pipeline achieves 544/544 detection at the block-or-escalate level. Because
InjecAgent's cases are single-step (a benign tool call whose output carries
an injected instruction, followed by the expected attacker action), this
result exercises VAULT and EXECUTOR's deterministic layer directly rather
than CHANNEL's multi-step provenance mechanism, which the two evaluations
above target specifically.

\subsection{Cost}
\label{sec:cost}
Deterministic parsing adds negligible cost on the hot path: the
\texttt{bashlex}-based AST parse used for \texttt{sh}/\texttt{bash}
(\S\ref{sec:vault-executor}) measures $\sim$0.5\,ms per call, and the
structural tokenizer used for other dialects is comparably cheap. Kavach
needs no GPU: the BGE embedding step used by the cosine-triage layer and by
NAVIGATOR runs CPU-only end-to-end, and every number in this section is
measured that way.
